\documentclass[a4paper,12pt]{article}

% ============================================================
% Кодировка, язык, математика
% ============================================================
\usepackage[T2A]{fontenc}
\usepackage[utf8]{inputenc}
\usepackage[russian]{babel}
\usepackage{amsmath,amssymb,amsthm,mathtools}
\usepackage{icomma}

% ============================================================
% Типографика
% ============================================================
\usepackage{helvet}
\renewcommand{\familydefault}{\sfdefault}
\usepackage{microtype}
\usepackage{setspace}
\onehalfspacing
\usepackage{parskip}

% ============================================================
% Страница
% ============================================================
\usepackage[
  a4paper,
  left=20mm,
  right=20mm,
  top=18mm,
  bottom=22mm
]{geometry}

% ============================================================
% Графика
% ============================================================
\usepackage{tikz}
\usetikzlibrary{calc}
\usetikzlibrary{arrows.meta,positioning,shapes.geometric}

\usepackage{tkz-euclide}

\usepackage{pgfplots}
\pgfplotsset{compat=1.18}

% ============================================================
% Таблицы
% ============================================================
\usepackage{tabularx,booktabs,longtable}
\renewcommand{\arraystretch}{1.35}

% ============================================================
% Списки, ссылки, колонтитулы
% ============================================================
\usepackage{enumitem}
\usepackage{hyperref}
\usepackage{fancyhdr}
\usepackage{xcolor}

% ============================================================
% Цветовая палитра
% ============================================================
\definecolor{accent}{HTML}{008080}
\definecolor{darkaccent}{HTML}{004D4D}
\definecolor{textgray}{HTML}{333333}
\definecolor{softgray}{HTML}{F2F5F5}
\definecolor{lightaccent}{HTML}{DCEEEE}

\color{textgray}

% ============================================================
% Hyperref
% ============================================================
\hypersetup{
  colorlinks=true,
  linkcolor=darkaccent,
  urlcolor=darkaccent
}

% ============================================================
% Колонтитулы
% ============================================================
\pagestyle{fancy}
\fancyhf{}
\fancyhead[L]{\small\color{textgray}Физика}
\fancyhead[R]{\small\color{textgray}Равноускоренное движение}
\fancyfoot[C]{\small\thepage}
\renewcommand{\headrulewidth}{0.3pt}
\renewcommand{\headrule}{\hbox to\headwidth{%
  \color{accent}\leaders\hrule height \headrulewidth\hfill}}

% ============================================================
% Списки
% ============================================================
\setlist[itemize]{
  leftmargin=7mm,
  itemsep=3pt,
  topsep=4pt
}

\setlist[enumerate]{
  leftmargin=8mm,
  itemsep=5pt,
  topsep=4pt
}

% ============================================================
% Заголовки
% ============================================================
\newcommand{\sectiontitle}[1]{%
  \vspace{2mm}
  {\Large\bfseries\color{darkaccent}#1}\par
  \vspace{1mm}
  {\color{accent}\rule{\linewidth}{0.5pt}}
  \vspace{3mm}
}

\newcommand{\subsectiontitle}[1]{%
  \vspace{2mm}
  {\large\bfseries\color{darkaccent}#1}\par
  \vspace{1mm}
}

\newcommand{\important}[1]{%
  {\bfseries\color{darkaccent}#1}
}

% ============================================================
% TikZ: блок-схемы
% ============================================================
\tikzset{
  flowstart/.style={
    ellipse,
    draw=accent,
    line width=0.9pt,
    minimum width=42mm,
    minimum height=10mm,
    align=center,
    inner sep=4pt
  },
  flowaction/.style={
    rectangle,
    rounded corners=3pt,
    draw=accent,
    line width=0.9pt,
    text width=104mm,
    minimum height=12mm,
    align=center,
    inner sep=5pt
  },
  flowdecision/.style={
    diamond,
    aspect=2.4,
    draw=accent,
    line width=0.9pt,
    text width=49mm,
    align=center,
    inner sep=2pt
  },
  flowarrow/.style={
    -{Stealth[length=3mm]},
    line width=0.9pt,
    draw=darkaccent
  }
}

% ============================================================
% Документ
% ============================================================
\begin{document}

% ============================================================
% ТИТУЛЬНЫЙ ЛИСТ
% ============================================================
\thispagestyle{empty}

\begin{center}

\vspace*{27mm}

{\Large\color{accent}\textbf{ФИЗИКА}}

\vspace{13mm}

{\Huge\bfseries\color{darkaccent}
Чек-лист}

\vspace{5mm}

{\LARGE\bfseries
Равноускоренное движение}

\vspace{4mm}

{\Large
Координата, путь, скорость и ускорение}

\vspace{20mm}

{\large
\textbf{Ключевые навыки занятия}}

\vspace{5mm}

\begin{minipage}{0.78\linewidth}
\centering
\large
читать уравнение движения;

\vspace{2mm}
находить \(x_0\), \(v_{0x}\) и \(a_x\);

\vspace{2mm}
составлять уравнение координаты;

\vspace{2mm}
выбирать формулу для пути по известным величинам;

\vspace{2mm}
переводить текст условия на язык физики.
\end{minipage}

\vfill

{\large \today}

\vspace{9mm}

{\large
Лёвин Артём Александрович\\[2mm]
\textbf{эксклюзивно для Володи Хачатуряна}}

\vspace{25mm}

\begin{tikzpicture}[remember picture,overlay]
  \draw[
    accent,
    line width=1.1pt
  ]
  ($(current page.south west)+(18mm,15mm)$)
  .. controls
  ($(current page.south west)+(65mm,27mm)$)
  and
  ($(current page.south east)+(-70mm,6mm)$)
  ..
  ($(current page.south east)+(-18mm,18mm)$);
\end{tikzpicture}

\end{center}

\clearpage

% ============================================================
% ОСНОВНОЙ ЧЕК-ЛИСТ
% ============================================================
\sectiontitle{1. Что означает ускорение}

Ускорение показывает, насколько быстро изменяется скорость тела:

\[
a_x=\frac{\Delta v_x}{\Delta t}.
\]

Для равноускоренного движения ускорение постоянно:

\[
a_x=\text{const}.
\]

\begin{itemize}
  \item \(a_x=0\) --- скорость вдоль оси \(Ox\) остаётся постоянной;
  \item \(a_x>0\) --- проекция скорости \(v_x\) со временем увеличивается;
  \item \(a_x<0\) --- проекция скорости \(v_x\) со временем уменьшается.
\end{itemize}

\important{Следите именно за проекциями.}
Знаки \(v_x\) и \(a_x\) зависят от выбранного направления оси \(Ox\).

\subsectiontitle{Изменение скорости}

При постоянном ускорении:

\[
\Delta v_x=a_x t.
\]

Следовательно,

\[
v_x=v_{0x}+a_x t.
\]

Здесь:
\[
v_{0x} \text{ --- начальная скорость},\qquad
v_x \text{ --- скорость через время }t.
\]

\sectiontitle{2. Путь и координата при равноускоренном движении}

Перемещение вдоль оси \(Ox\):

\[
\boxed{
s_x=v_{0x}t+\frac{a_xt^2}{2}
}
\]

Координата связана с перемещением формулой

\[
x=x_0+s_x.
\]

Поэтому основное уравнение координаты имеет вид

\[
\boxed{
x=x_0+v_{0x}t+\frac{a_xt^2}{2}
}
\]

где

\[
x_0 \text{ --- начальная координата}.
\]

\important{Единицы СИ:}

\[
[x]=[s]=\text{м},\qquad
[v]=\frac{\text{м}}{\text{с}},\qquad
[a]=\frac{\text{м}}{\text{с}^2},\qquad
[t]=\text{с}.
\]

\subsectiontitle{Пример}

Пусть

\[
x=-2+0{,}5t+10t^2.
\]

Сравниваем с

\[
x=x_0+v_{0x}t+\frac{a_xt^2}{2}.
\]

Получаем:

\[
x_0=-2\text{ м},
\]

\[
v_{0x}=0{,}5\frac{\text{м}}{\text{с}},
\]

\[
\frac{a_x}{2}=10.
\]

Следовательно,

\[
\boxed{
a_x=20\frac{\text{м}}{\text{с}^2}
}
\]

\important{Ключевой момент:}
коэффициент перед \(t^2\) равен

\[
\frac{a_x}{2}.
\]

Поэтому ускорение находится удвоением этого коэффициента.

\clearpage

% ============================================================
% БЛОК-СХЕМА 1
% ============================================================
\thispagestyle{plain}

\begin{center}
{\LARGE\bfseries\color{darkaccent}
Алгоритм: как прочитать уравнение координаты}
\end{center}

\vspace{7mm}

\begin{center}
\begin{tikzpicture}[
  node distance=9mm,
  every node/.style={font=\small}
]

\node[flowstart] (start)
{Дано уравнение \(x(t)\)};

\node[flowaction,below=of start] (standard)
{Запишите стандартный вид:
\[
x=x_0+v_{0x}t+\frac{a_xt^2}{2}
\]};

\node[flowaction,below=of standard] (constant)
{Свободный член --- это \(x_0\)};

\node[flowaction,below=of constant] (linear)
{Коэффициент перед \(t\) --- это \(v_{0x}\)};

\node[flowaction,below=of linear] (square)
{Коэффициент перед \(t^2\) --- это \(\dfrac{a_x}{2}\)};

\node[flowaction,below=of square] (acc)
{Умножьте коэффициент при \(t^2\) на \(2\):
\[
a_x=2\cdot(\text{коэффициент при }t^2)
\]};

\node[flowstart,below=of acc] (finish)
{Запишите \(x_0\), \(v_{0x}\), \(a_x\) с единицами};

\draw[flowarrow] (start) -- (standard);
\draw[flowarrow] (standard) -- (constant);
\draw[flowarrow] (constant) -- (linear);
\draw[flowarrow] (linear) -- (square);
\draw[flowarrow] (square) -- (acc);
\draw[flowarrow] (acc) -- (finish);

\end{tikzpicture}
\end{center}

\clearpage

% ============================================================
% ПРОДОЛЖЕНИЕ ЧЕК-ЛИСТА
% ============================================================
\sectiontitle{3. Обратная задача: составить уравнение координаты}

Если известны

\[
x_0,\qquad v_{0x},\qquad a_x,
\]

используйте

\[
x=x_0+v_{0x}t+\frac{a_xt^2}{2}.
\]

\subsectiontitle{Пример}

Дано:

\[
x_0=10\text{ м},
\qquad
v_{0x}=-2\frac{\text{м}}{\text{с}},
\qquad
a_x=3\frac{\text{м}}{\text{с}^2}.
\]

Подставляем:

\[
x=10-2t+\frac{3t^2}{2}.
\]

Ответ:

\[
\boxed{
x(t)=10-2t+\frac{3}{2}t^2
}
\]

\sectiontitle{4. Как найти координату в заданный момент}

Используйте

\[
x=x_0+v_{0x}t+\frac{a_xt^2}{2}.
\]

\subsectiontitle{Пример}

Пусть

\[
x_0=-10\text{ м},\qquad
v_{0x}=1\frac{\text{м}}{\text{с}},\qquad
a_x=1\frac{\text{м}}{\text{с}^2},
\qquad
t=10\text{ с}.
\]

Тогда

\[
x=-10+1\cdot10+\frac{1\cdot10^2}{2}.
\]

\[
x=-10+10+50=50\text{ м}.
\]

\[
\boxed{x=50\text{ м}}
\]

\sectiontitle{5. Как найти время по заданной координате}

Если известна координата, в которой должно оказаться тело, подставьте её вместо \(x\):

\[
x=x_0+v_{0x}t+\frac{a_xt^2}{2}.
\]

После подстановки часто получается квадратное уравнение относительно \(t\).

Например:

\[
100=10-2t+\frac{3t^2}{2}.
\]

Переносим всё в одну часть:

\[
\frac{3}{2}t^2-2t-90=0.
\]

После устранения дроби:

\[
3t^2-4t-180=0.
\]

Для уравнения

\[
At^2+Bt+C=0
\]

используем

\[
D=B^2-4AC,
\]

\[
t_{1,2}=
\frac{-B\pm\sqrt D}{2A}.
\]

В физических задачах дополнительно проверяем смысл корней.

Обычно рассматривается

\[
\boxed{t\ge0}.
\]

Если один из корней отрицательный и задача описывает движение после момента \(t=0\), такой корень не соответствует рассматриваемому промежутку движения.

\clearpage

% ============================================================
% БЛОК-СХЕМА 2
% ============================================================
\thispagestyle{plain}

\begin{center}
{\LARGE\bfseries\color{darkaccent}
Алгоритм: координата задана, требуется время}
\end{center}

\vspace{6mm}

\begin{center}
\begin{tikzpicture}[
  node distance=8mm,
  every node/.style={font=\small}
]

\node[flowstart] (start)
{Известна требуемая координата \(x\)};

\node[flowaction,below=of start] (formula)
{Запишите:
\[
x=x_0+v_{0x}t+\frac{a_xt^2}{2}
\]};

\node[flowaction,below=of formula] (sub)
{Подставьте известные величины};

\node[flowaction,below=of sub] (eq)
{Приведите уравнение к виду
\[
At^2+Bt+C=0
\]};

\node[flowdecision,below=11mm of eq] (disc)
{\(D\ge0\)?};

\node[flowaction,below=12mm of disc] (roots)
{Найдите корни:
\[
t_{1,2}=\frac{-B\pm\sqrt D}{2A}
\]};

\node[flowaction,below=of roots] (physics)
{Проверьте физический смысл:
\(t\ge0\) и соответствие условию};

\node[flowstart,below=of physics] (finish)
{Запишите допустимое время};

\node[flowstart,right=25mm of disc] (noroot)
{Решения для данного движения нет};

\draw[flowarrow] (start) -- (formula);
\draw[flowarrow] (formula) -- (sub);
\draw[flowarrow] (sub) -- (eq);
\draw[flowarrow] (eq) -- (disc);

\draw[flowarrow] (disc) -- node[right,font=\small]{да} (roots);
\draw[flowarrow] (roots) -- (physics);
\draw[flowarrow] (physics) -- (finish);

\draw[flowarrow]
  (disc.east)
  -- ++(8mm,0)
  |- node[pos=0.20,above,font=\small]{нет} (noroot.west);

\end{tikzpicture}
\end{center}

\clearpage

% ============================================================
% ФОРМУЛЫ ДЛЯ ПУТИ
% ============================================================
\sectiontitle{6. Три основные формулы для перемещения}

Для равноускоренного движения полезны три формы одной и той же кинематической связи.

\subsectiontitle{Формула 1. Известны \(v_{0x}\), \(a_x\), \(t\)}

\[
\boxed{
s_x=v_{0x}t+\frac{a_xt^2}{2}
}
\]

Здесь отсутствует конечная скорость \(v_x\).

\subsectiontitle{Формула 2. Известны \(v_{0x}\), \(v_x\), \(a_x\)}

\[
\boxed{
s_x=\frac{v_x^2-v_{0x}^2}{2a_x}
}
\]

Формула применима при

\[
a_x\ne0.
\]

Здесь время \(t\) для вычисления перемещения не требуется.

\subsectiontitle{Формула 3. Известны \(v_{0x}\), \(v_x\), \(t\)}

Средняя скорость при равноускоренном движении:

\[
v_{\text{ср},x}
=
\frac{v_{0x}+v_x}{2}.
\]

Поэтому

\[
\boxed{
s_x=
\frac{v_{0x}+v_x}{2}\,t
}
\]

Здесь ускорение для расчёта перемещения не требуется.

\subsectiontitle{Как выбирать}

\begin{center}
\begin{tabularx}{\linewidth}{@{}X X@{}}
\toprule
\textbf{Известно} & \textbf{Удобная формула} \\
\midrule

\(v_{0x},\,a_x,\,t\)
&
\(\displaystyle
s_x=v_{0x}t+\frac{a_xt^2}{2}
\)
\\[3mm]

\(v_{0x},\,v_x,\,a_x\)
&
\(\displaystyle
s_x=\frac{v_x^2-v_{0x}^2}{2a_x}
\)
\\[3mm]

\(v_{0x},\,v_x,\,t\)
&
\(\displaystyle
s_x=\frac{v_{0x}+v_x}{2}\,t
\)
\\

\bottomrule
\end{tabularx}
\end{center}

\important{Полезный принцип:}
сначала посмотрите, какие величины даны и какую требуется найти. Затем выбирайте формулу, содержащую минимум неизвестных.

\clearpage



% ============================================================
% ТЕКСТОВЫЕ ЗАДАЧИ
% ============================================================
\sectiontitle{7. Перевод текста задачи на язык физики}

Один из важнейших навыков --- превращать каждую содержательную фразу условия в физическую запись.

\begin{center}
\begin{tabularx}{\linewidth}{@{}X X@{}}
\toprule
\textbf{Фраза в условии} & \textbf{Физическая запись} \\
\midrule

«начинает движение из состояния покоя»
&
\[
v_{0x}=0
\]
\\

«остановилось»
&
\[
v_x=0
\]
в соответствующий момент времени
\\

«движется равноускоренно»
&
\[
a_x=\text{const}
\]
\\

«скорость увеличилась на \(30\text{ м/с}\)»
&
\[
\Delta v_x=30\frac{\text{м}}{\text{с}}
\]
\\

«через \(6\) секунд»
&
\[
t=6\text{ с}
\]
\\

\bottomrule
\end{tabularx}
\end{center}

\subsectiontitle{Пример: старт из состояния покоя}

Лыжник начинает спуск из состояния покоя. Через \(30\) с его скорость становится равной \(15\text{ м/с}\).

Записываем:

\[
v_{0x}=0,
\qquad
v_x=15\frac{\text{м}}{\text{с}},
\qquad
t=30\text{ с}.
\]

Используем

\[
a_x=\frac{v_x-v_{0x}}{t}.
\]

Тогда

\[
a_x=
\frac{15-0}{30}
=
0{,}5\frac{\text{м}}{\text{с}^2}.
\]

\[
\boxed{
a_x=0{,}5\frac{\text{м}}{\text{с}^2}
}
\]

\subsectiontitle{Пример: увеличение скорости}

Тело движется с постоянным ускорением

\[
a_x=5\frac{\text{м}}{\text{с}^2}.
\]

Найдём, на сколько изменится его скорость за \(6\) с.

Используем

\[
\Delta v_x=a_xt.
\]

\[
\Delta v_x=5\cdot6=30\frac{\text{м}}{\text{с}}.
\]

\[
\boxed{
\Delta v_x=30\frac{\text{м}}{\text{с}}
}
\]

Здесь найдено именно \important{изменение скорости}. Для нахождения конечной скорости дополнительно потребовалась бы начальная скорость:

\[
v_x=v_{0x}+\Delta v_x.
\]

\sectiontitle{8. Контроль здравого смысла}

Перед окончательной записью ответа задайте себе несколько вопросов.

\begin{itemize}
  \item Согласуются ли знаки скорости и ускорения с описанием движения?
  \item Если скорость по выбранной оси уменьшается от \(40\) до \(30\text{ м/с}\), соответствует ли этому знак ускорения?
  \item Получилось ли положительное время для процесса, происходящего после \(t=0\)?
  \item Сохранились ли точные дроби до конца вычислений?
  \item Совпадают ли единицы ответа с искомой величиной?
\end{itemize}

Например, если

\[
v_{0x}=40\frac{\text{м}}{\text{с}},
\qquad
v_x=30\frac{\text{м}}{\text{с}},
\qquad
a_x=-3\frac{\text{м}}{\text{с}^2},
\]

то уменьшение скорости согласуется с отрицательным ускорением.

Перемещение удобно найти сразу:

\[
s_x=
\frac{v_x^2-v_{0x}^2}{2a_x}.
\]

\[
s_x=
\frac{30^2-40^2}{2(-3)}
=
\frac{900-1600}{-6}
=
\frac{-700}{-6}
=
\frac{350}{3}\text{ м}.
\]

\[
\boxed{
s_x=\frac{350}{3}\text{ м}
}
\]

Сохранение обыкновенной дроби позволяет избежать накопления ошибок округления.

\sectiontitle{9. Типичные ошибки}

\begin{enumerate}
  \item Принимать коэффициент при \(t^2\) за ускорение:
  \[
  x=x_0+v_{0x}t+\frac{a_xt^2}{2}.
  \]
  Коэффициент перед \(t^2\) равен \(\dfrac{a_x}{2}\).

  \item Терять знак начальной скорости или ускорения.

  \item Подменять изменение скорости конечной скоростью:
  \[
  \Delta v_x=v_x-v_{0x}.
  \]

  \item Сразу переводить точные дроби в десятичные и преждевременно округлять.

  \item Искать промежуточные величины, когда существует более короткая формула.

  \item Использовать формулу
  \[
  s_x=\frac{v_x^2-v_{0x}^2}{2a_x}
  \]
  при \(a_x=0\). В таком случае знаменатель обращается в ноль.

  \item Сохранять отрицательное значение времени без проверки физического смысла.
\end{enumerate}

\sectiontitle{10. Минимум, который нужно знать наизусть}

\[
\boxed{
v_x=v_{0x}+a_xt
}
\]

\[
\boxed{
s_x=v_{0x}t+\frac{a_xt^2}{2}
}
\]

\[
\boxed{
x=x_0+v_{0x}t+\frac{a_xt^2}{2}
}
\]

\[
\boxed{
s_x=\frac{v_x^2-v_{0x}^2}{2a_x}
}
\qquad (a_x\ne0)
\]

\[
\boxed{
s_x=\frac{v_{0x}+v_x}{2}\,t
}
\]

\[
\boxed{
a_x=\frac{\Delta v_x}{\Delta t}
}
\]

\[
\boxed{
\Delta v_x=v_x-v_{0x}
}
\]

\clearpage

% ============================================================
% ТРЕНИРОВОЧНЫЙ БЛОК
% ============================================================
\thispagestyle{plain}

\begin{center}
{\LARGE\bfseries\color{darkaccent}
Тренировочный блок}

\vspace{2mm}

{\large
Путь А --- первичное закрепление}
\end{center}

\vspace{6mm}

\textbf{Задание 1. Чтение уравнения движения}

Координата тела изменяется по закону

\[
x=-4+6t-3t^2,
\]

где \(x\) измеряется в метрах, \(t\) --- в секундах.

Определите:

\begin{enumerate}[label=\alph*)]
  \item начальную координату \(x_0\);
  \item начальную скорость \(v_{0x}\);
  \item ускорение \(a_x\).
\end{enumerate}

\vspace{7mm}

\textbf{Задание 2. Координата тела}

Тело движется равноускоренно вдоль оси \(Ox\).

Известно:

\[
x_0=-8\text{ м},
\qquad
v_{0x}=4\frac{\text{м}}{\text{с}},
\qquad
a_x=2\frac{\text{м}}{\text{с}^2}.
\]

Найдите координату тела через

\[
t=5\text{ с}.
\]

\vspace{7mm}

\textbf{Задание 3. Выбор рациональной формулы}

Автомобиль движется прямолинейно равноускоренно.

Его скорость уменьшилась от

\[
v_{0x}=25\frac{\text{м}}{\text{с}}
\]

до

\[
v_x=15\frac{\text{м}}{\text{с}}.
\]

Ускорение автомобиля:

\[
a_x=-2\frac{\text{м}}{\text{с}^2}.
\]

Найдите перемещение автомобиля за этот промежуток движения.

Выберите формулу, позволяющую выполнить расчёт без предварительного нахождения времени.

\clearpage

% ============================================================
% ОТВЕТЫ
% ============================================================
\thispagestyle{plain}

\begin{center}
{\LARGE\bfseries\color{darkaccent}
Ответы к тренировочному блоку}
\end{center}

\vspace{9mm}

\begin{table}[ht]
\centering
\begin{tabularx}{\linewidth}{@{}c X@{}}
\toprule
\textbf{№} & \textbf{Ответ} \\
\midrule

1 &
\[
x_0=-4\text{ м},
\qquad
v_{0x}=6\frac{\text{м}}{\text{с}},
\qquad
a_x=-6\frac{\text{м}}{\text{с}^2}.
\]
\\[5mm]

2 &
\[
x=37\text{ м}.
\]
\\[5mm]

3 &
\[
s_x=100\text{ м}.
\]
\\

\bottomrule
\end{tabularx}
\end{table}

\vfill

\begin{center}
\large\color{darkaccent}
\textbf{Для следующего занятия}

\vspace{4mm}

Выучить основные формулы равноускоренного движения
и потренироваться выбирать формулу по набору известных величин.

\vspace{4mm}

Особое внимание:

\[
\boxed{
\text{коэффициент при }t^2=\frac{a_x}{2}
}
\]

и

\[
\boxed{
\Delta v_x=v_x-v_{0x}
}
\]

\end{center}

\end{document}